% Page 006 — page imprimée 2
% Reconstruction à partir du scan.
% Cette page est la deuxième page numérotée du livre.


Mais il n’était pas question de se limiter à cette période du 17\textsuperscript{o}-18\textsuperscript{o} siècle, marquée par l’arrivée de la Réforme, ses progrès rapides suivis par des mesures de restriction des libertés de plus en plus draconiennes à Meyrueis même, aboutissant à La Révocation de l’Edit de Nantes en 1685 avec la destruction du temple et les abjurations forcées de la plupart des meyrueisiens subissant les dragonnades.

Arrive la période des philosophes et des Lumières qui débouchera sur la Révolution, vécue comme une libération par les protestants du sud de la Lozère. Ne créèrent-ils pas à Meyrueis une « Société Populaire » dès 1793, où se font déjà jour, d’ailleurs, les luttes entre Girondins et Montagnards ? Révolution combattue par les Catholiques, dans la haute Lozère et les Gorges du Tarn, surtout à la suite de la Constitution Civile du Clergé. Il y eut bien des excès de part et d’autre. Mais aux morts républicains répondirent de nombreuses arrestations et condamnations de royalistes. La Malène en souffrit particulièrement car nombre d'hommes de ce village des gorges du Tarn s’étaient engagés dans les bandes royalistes entraînées par Charrier ancien député de la Constituante et 37 furent fusillés à Florac. Il y est fait allusion dans les pages qui suivent.

Vous ne trouverez pas grand chose sur le 19\textsuperscript{o} siècle. Il faudrait y revenir…

Enfin nous arrivons à la période contemporaine, celle que nous avons connue ou que nos pères ont connue au cours des années 1900. Elles commencent avec l'Affaire Dreyfus et les décrets de séparation des Eglises et de l’Etat en 1905.

Des anciens nous ont parlé ou même ont laissé par écrit, des souvenirs précieux, relatés dans cette brochure qui n’est que l’ébauche de ce que pourrait être un ouvrage sinon exhaustif, du moins plus complet, mieux documenté et présentant une véritable iconographie, Nous appelons de nos vœux, une véritable et complète histoire de Meyrueis présentant la vie et les activités de « Meyrueis, moun villetto » depuis les origines jusqu’à nos jours.

\vspace{0.35em}

\begin{tcolorbox}[
  colback=white,
  colframe=black,
  boxrule=0.7pt,
  arc=0pt,
  left=0.25cm,
  right=0.25cm,
  top=0.12cm,
  bottom=0.12cm,
  boxsep=0pt
]
\small
Les pages qui suivent ont été rédigées au cours des quatre ou cinq dernières années, au fur et à mesure, du dépouillement de documents consultés aux Archives de Mende, ou de Montpellier. D’autres renseignements ont été recueillis dans plusieurs livres et publications à la Bibliothèque de la Société d’Histoire du Protestantisme Français (SHPF) à Paris ainsi qu’à la Médiathèque de Montpellier. Enfin sont présentés quelques documents inédits.

C’est la raison pour laquelle vous trouverez des répétitions, et des renseignements qui ne se suivent pas toujours de façon chronologique.
\end{tcolorbox}

